\chapter{Programs — Copy}
%%% generated by the blueprint pipeline from TCSlib.Programs.Light.Lib.Copy
\section{Overview}

Declarations of the Lean module \texttt{TCSlib.Programs.Light.Lib.Copy}.

\section{Declarations}

\begin{definition}[Source index]
\lean{Light.Copy.Src}
\label{Light.Copy.Src}
\leanok
The source index is the natural number $0$.
\end{definition}

\begin{definition}[Copy destination index]
\lean{Light.Copy.Dst}
\label{Light.Copy.Dst}
\leanok
\uses{Light.Copy.Src}
The copy destination index is the natural number $1$.
\end{definition}

\begin{definition}[Copy length]
\lean{Light.Copy.Len}
\label{Light.Copy.Len}
\leanok
\uses{Light.Copy.Src}
The copy length is the natural number $2$.
\end{definition}

\begin{definition}[Copy source index]
\lean{Light.Copy.Idx}
\label{Light.Copy.Idx}
\leanok
\uses{Light.Copy.Src}
The copy source index is the natural number $3$.
\end{definition}

\begin{definition}[Copy body]
\lean{Light.copyBody}
\label{Light.copyBody}
\leanok
\uses{Light.Copy.Dst, Light.Copy.Idx, Light.Copy.Len, Light.Copy.Src, Light.M, Light.Stmt, Light.pass, Light.v}
The copy body is the statement obtained by applying the pass construction with index
$3$, length $2$, destination index $1$, and source address $0+3$.
\end{definition}

\begin{definition}[Copy running time]
\lean{Light.copyTime}
\label{Light.copyTime}
\leanok
For each natural number $n$, the copy running time is $16n+6$.
\end{definition}

\begin{theorem}[Bounded-time memory copy]
\lean{Light.copy_meets}
\label{Light.copy_meets}
\leanok
\uses{Light.Apart, Light.Ends.pass, Light.Meets, Light.Meets.of_body, Light.SameOutside, Light.copyBody, Light.copyTime, Light.wrote, Light.wrote_done, Light.wrote_rest}
\difficulty{3}
Let $P$ be a program, let $p,d\in\mathbb{N}$, and let $\mathrm{lim}$ specify word,
space, and depth limits. Suppose the program entry at $p$ is the copy routine, the space
limit is at most the word limit, and the source interval $[\mathrm{src},\mathrm{src}+n)$
and destination interval $[\mathrm{dst},\mathrm{dst}+n)$ lie within the space bound and
are disjoint. Then, starting with arguments $(\mathrm{src},\mathrm{dst},n)$ and memory
$\mu$, that routine has an execution at depth $d$ of cost at most $16n+6$ whose final
memory agrees with $\mu$ outside the destination interval and satisfies
$\mu'(\mathrm{dst}+i)=\mu(\mathrm{src}+i)$ for every $i<n$.
\end{theorem}

\begin{definition}[Fill destination index]
\lean{Light.Fill.Dst}
\label{Light.Fill.Dst}
\leanok
The fill destination index is the natural number $0$.
\end{definition}

\begin{definition}[Fill length]
\lean{Light.Fill.Len}
\label{Light.Fill.Len}
\leanok
\uses{Light.Fill.Dst}
The fill length is the natural number $1$.
\end{definition}

\begin{definition}[Fill value]
\lean{Light.Fill.Val}
\label{Light.Fill.Val}
\leanok
\uses{Light.Fill.Dst}
The fill value is the natural number $2$.
\end{definition}

\begin{definition}[Index constant]
\lean{Light.Fill.Idx}
\label{Light.Fill.Idx}
\leanok
\uses{Light.Fill.Dst}
The index is the natural number $3$.
\end{definition}

\begin{definition}[Fill body construction]
\lean{Light.fillBody}
\label{Light.fillBody}
\leanok
\uses{Light.Fill.Dst, Light.Fill.Idx, Light.Fill.Len, Light.Fill.Val, Light.Stmt, Light.pass, Light.v}
The fill body is the statement obtained by applying the pass routine to the index and to
the variable expressions for length, destination, and value.
\end{definition}

\begin{definition}[Fill time]
\lean{Light.fillTime}
\label{Light.fillTime}
\leanok
For each natural number $n$, the fill time is $13n+6$.
\end{definition}

\begin{theorem}[Filling a Memory Segment]
\lean{Light.fill_meets}
\label{Light.fill_meets}
\leanok
\uses{Light.Ends.pass, Light.Meets, Light.Meets.of_body, Light.SameOutside, Light.Seg, Light.fillBody, Light.fillTime, Light.seg_wrote}
\difficulty{1}
Let a program with limits $L$ have procedure $p$ equal to the specified fill procedure.
For any initial memory $\mu:\mathbb N\to\mathbb Z$, destination $a\in\mathbb N$, length
$n\in\mathbb N$, and value $x\in\mathbb Z$, suppose the word bound is at least the space
bound and $a+n\le L.space$. On input $(a,n,x)$, the procedure has an execution taking at
most $13n+6$ steps whose final memory contains $n$ consecutive copies of $x$ starting at
$a$ and agrees with $\mu$ at every location outside those $n$ positions.
\end{theorem}
